Some goblins, $N$ in number, are standing in a row while ``trick-or-treat"-ing.
Each goblin is at all times either $2'$ tall or $3'$ tall, but can change
spontaneously from one of these two heights to the other at will. While lined
up in such a row, a goblin is called a \textbf{Local Giant Goblin} (LGG) if
he/she/it is not standing beside a taller goblin. Let $G(N)$ be the total of
all occurrences of LGG's as the row of $N$ goblins transmogrifies through all
possible distinct configurations, where height is the only distinguishing
characteristic. As an example, with $N = 2$, the distinct configurations are
$\hat{2}\hat{2}$, $2\hat{3}$, $\hat{3}2$, $\hat{3}\hat{3}$, where a cap
indicates an LGG. Thus $G(2) = 6$.

\begin{enumerate}[(a)]
\item Find $G(3)$ and $G(4)$.
\item Find, with proof, the general formula for $G(N)$, $N = 1, 2, 3, \dots$.
\end{enumerate}
